
\subsection{算子 \texorpdfstring{$\mathfrak p$}{p} 的推广}
\label{subsec:3-8-9}
\liftIndex{I355}

把算子 $\mathfrak p_k$ 与 $\mathfrak q$ 结合，可以得到算子
$\mathfrak p_{\ell,k}$，从而推广 $\mathfrak p_k$。（算子
$\mathfrak p$ 与 $\mathfrak q$ 的区别如下：定义 $\mathfrak p$ 所用的
模空间以第 $0$ 个标记点为内部标记点；而定义 $\mathfrak q$ 所用的
模空间以第 $0$ 个标记点为边界标记点。）第 6 章将用算子
$\mathfrak p_{\ell,k}$ 描述那里构造的谱序列微分之像。

由于 $\mathfrak p_{\ell,k}$ 的构造完全类似于本节已经讨论过的构造，
这里仅写出结论。先再回顾一点格罗莫夫--威腾不变量。令
$\mathcal M_{0,\ell+2}(M;\widetilde\beta)$ 为同调类
$\widetilde\beta\in H_2(M)$、带有 $\ell+2$ 个标记点的伪全纯球面
模空间。存在自然的评价映射
\[
 ev=(ev_1,\ldots,ev_{\ell+1},ev_0)\colon
 \mathcal M_{0,\ell+2}(M;\widetilde\beta)
 \longrightarrow M^{\ell+2}.
\]
对 $M$ 中的链 $Q_i$（$i=1,\ldots,\ell+1$），置
\[
\begin{aligned}
 &\mathcal M_{0,\ell+2}
 (M;\widetilde\beta;\vec Q)\\
 &\quad=
 \mathcal M_{0,\ell+2}(M;\widetilde\beta)
 _{(ev_1,\ldots,ev_{\ell+1})}
 \mathbin{\times_{M^{\ell+1}}}
 (Q_1\times\cdots\times Q_{\ell+1}),
\end{aligned}
\]
并定义
\begin{equation}
\begin{aligned}
 GW_{0,\ell+1}(M)(\vec Q)
 ={}&\sum_{\widetilde\beta}
 T^{\omega[\widetilde\beta]}
 e^{c_1(M)(\widetilde\beta)}\\
 &\quad\cdot ev_{0*}\bigl(
 \mathcal M_{0,\ell+2}
 (M;\widetilde\beta;\vec Q)
 \bigr).
\end{aligned}
\tag{3.8.77}\label{eq:3-8-77}
\end{equation}
式 \eqref{eq:3-8-77} 推广了 \eqref{eq:3-8-2}。（取适当的三角剖分后，
把右端视为 $S^*(M;\Lambda_{0,\mathrm{nov}})$ 的元素。）

还要指出，正如 $\mathfrak m_k$ 在引理--定义 3.8.4 中诱导
$\widehat d^{\mathrm{cyc}}$，$\mathfrak q$ 也诱导
\[
\begin{aligned}
 \widehat{\mathfrak q}^{\mathrm{cyc}}\colon{}&
 \widehat EC(M;\Lambda_{0,\mathrm{nov}})[2]
 \otimes
 B^{\mathrm{cyc}}C_{1+}(L;\Lambda_{0,\mathrm{nov}})[1]\\
 &\longrightarrow
 B^{\mathrm{cyc}}C_{1+}(L;\Lambda_{0,\mathrm{nov}})[1].
\end{aligned}
\]

\begin{proposition}{3.8.78}\label{prop:3-8-78}
当 $R=\mathbb Q$ 时，存在次数为 $n+1$ 的一列算子
\[
\begin{aligned}
 \mathfrak p_{\ell,k}\colon{}&
 E_\ell C(M;\Lambda_{0,\mathrm{nov}})[2]
 \otimes
 B_k^{\mathrm{cyc}}C_{1+}(L;\Lambda_{0,\mathrm{nov}})[1]\\
 &\longrightarrow S^*(M;\Lambda_{0,\mathrm{nov}}),
 \qquad k,\ell=0,1,2,\ldots,
\end{aligned}
\]
具有下列性质：

\begin{align}
 \mathfrak p_{0,k}&=\mathfrak p_k,
 \tag{3.8.79.1}\label{eq:3-8-79-1}\\
 \mathfrak p_{\ell,k}&\equiv0
 \pmod{\Lambda^+_{0,\mathrm{nov}}}
 \qquad(\ell>0).
 \tag{3.8.79.2}\label{eq:3-8-79-2}
\end{align}

当 $k>0$ 时，有
\begin{equation}
\begin{aligned}
 &\mathfrak p_{\ell,k}
 \bigl(\widehat\delta_M(y)\otimes x\bigr)
 +\delta_M\bigl(\mathfrak p_{\ell,k}(y\otimes x)\bigr)\\
 &\quad+
 \sum_{c_1}
 \mathfrak p\left(
 y_{c_1}^{(2;1)}\otimes
 \widehat{\mathfrak q}^{\mathrm{cyc}}
 \bigl(y_{c_1}^{(2;2)}\otimes x\bigr)
 \right)=0.
\end{aligned}
\tag{3.8.79.3}\label{eq:3-8-79-3}
\end{equation}
这里的记号与定理 3.8.32 相同。此外，
\begin{equation}
\begin{aligned}
 &\mathfrak p_{\ell,0}
 \bigl(\widehat\delta_M(y)\otimes1\bigr)
 +\delta_M\bigl(\mathfrak p_{\ell,0}(y\otimes1)\bigr)\\
 &\quad+
 \sum_{c_1}
 (-1)^{\deg y_{c_1}^{(2;1)}}
 \mathfrak p\left(
 y_{c_1}^{(2;1)}\otimes
 \mathfrak q\bigl(y_{c_1}^{(2;2)}\otimes1\bigr)
 \right)\\
 &\quad+
 \frac1{\ell!}GW_{0,\ell+1}(M)(y\otimes L)=0.
\end{aligned}
\tag{3.8.79.4}\label{eq:3-8-79-4}
\end{equation}
左端最后一项按 \eqref{eq:3-8-77} 定义。具体地，若
$y=y_1\otimes\cdots\otimes y_\ell$，就取
$Q_i=y_i$（$i=1,\ldots,\ell$）以及 $Q_{\ell+1}=L$。

还可把 $\mathfrak p_{\ell,k}$ 延拓为
\[
\begin{aligned}
 \mathfrak p^+_{\ell,k}\colon{}&
 E_\ell C(M;\Lambda_{0,\mathrm{nov}})[2]
 \otimes
 B_k^{\mathrm{cyc}}C_{1+}(L;\Lambda_{0,\mathrm{nov}})^+[1]\\
 &\longrightarrow S^*(M;\Lambda_{0,\mathrm{nov}}),
\end{aligned}
\]
使
\begin{equation}
 \mathfrak p^+_{0,k}=\mathfrak p^+_k,
\tag{3.8.79.5}\label{eq:3-8-79-5}
\end{equation}
并满足与 (3.8.35.1)--(3.8.35.3) 类似的相容条件。特别地，
\begin{equation}
 \mathfrak p^+_{\ell,k}
 \bigl(y\otimes[\mathbf e^+,x_1,\ldots,x_{k-1}]\bigr)=0
 \qquad(\ell\ne0,\quad k\geq0).
\tag{3.8.79.6}\label{eq:3-8-79-6}
\end{equation}
这里与第 3.3 节一样，
\[
 C_{1+}(L;\Lambda_{0,\mathrm{nov}})^+
 =C_{1+}(L;\Lambda_{0,\mathrm{nov}})
 \oplus\Lambda_{0,\mathrm{nov}}\mathbf e^+
 \oplus\Lambda_{0,\mathrm{nov}}\mathbf f.
\]
\end{proposition}

算子 $\mathfrak p_{\ell,k}$ 的构造使用模空间
$\mathcal M_{k,\ell+1}^{\mathrm{main}}(L,\beta)$，其中最后一个（第 $0$
个）标记点是内部标记点。（回忆：定义 $\mathfrak p$ 时的最后一个标记点
是内部标记点，而定义 $\mathfrak q$ 时的最后一个标记点是边界标记点。）
见图 3.8.6。例如，图 3.8.7 给出
\[
\begin{aligned}
 \mathfrak p\Bigl(
 (Q_3\otimes Q_4)\otimes
 \bigl(&\mathfrak q(
 (Q_1\otimes Q_2)\otimes(P_6\otimes P_1\otimes P_2))\\
 &\otimes P_3\otimes P_4\otimes P_5
 \bigr)
 \Bigr),
\end{aligned}
\]
它是 \eqref{eq:3-8-79-3} 左端最后一项中的一个求和项\translatornote{原书式
(3.8.79.3) 的求和号写成
$\sum_{c_1,c_2}$，但被求和项只含 $c_1$，而本节约定中 $c_2$ 专门标记
第二个张量的余积分解；本式没有对 $x$ 作这类分解。紧随其后的
(3.8.79.4) 以及 $L_\infty$ 模同态的重述都只需对 $y$ 的余积分解求和，
译文据此删去未被使用的 $c_2$。}。

\begin{figure}[!htbp]
  \centering
  \setlength{\unitlength}{0.7mm}
  \begin{picture}(118,75)(0,-8)
    \put(59,36){\circle{48}}
    \put(55,60){\line(0,-1){3}}
    \put(75,53){\line(-1,-1){2.5}}
    \put(80,23){\line(-1,1){2.5}}
    \put(59,12){\line(0,1){3}}
    \put(38,20){\line(1,1){2.5}}
    \put(36,45){\line(1,-1){2.5}}
    \put(50,62){$P_3$}
    \put(77,54){$P_2$}
    \put(82,18){$P_1$}
    \put(55,5){$P_k$}
    \put(28,15){$P_5$}
    \put(25,46){$P_4$}
    \put(47,35){\circle*{1.7}}
    \put(59,43){\circle*{1.7}}
    \put(63,25){\circle*{1.7}}
    \put(47,36){\line(0,1){9}}
    \put(59,44){\line(1,2){4}}
    \put(63,26){\line(0,1){8}}
    \qbezier[8](47,34)(47,29)(47,26)
    \qbezier[8](59,42)(59,37)(59,34)
    \qbezier[8](63,24)(63,19)(63,16)
    \put(43,47){$Q_1$}
    \put(64,51){$Q_2$}
    \put(65,34){$Q_\ell$}
    \put(74,36){\circle*{1.7}}
    \put(77,36){\vector(1,0){7}}
    \put(77,40){$ev_0$}
    \qbezier[20](35,13)(43,3)(55,1)
  \end{picture}
  \[
   \mathfrak p_{\ell,k}
   \bigl((Q_1\otimes\cdots\otimes Q_\ell)
   \otimes(P_1\otimes\cdots\otimes P_k)\bigr)
  \]
  \renewcommand{\thefigure}{3.8.6}
  \refstepcounter{figure}
  \textbf{\figurename~\thefigure}\label{fig:3-8-6}
\end{figure}

\begin{figure}[!htbp]
  \centering
  \setlength{\unitlength}{0.75mm}
  \begin{picture}(126,62)(0,-3)
    \put(42,30){\circle{42}}
    \put(84,30){\circle{42}}
    \put(63,30){\circle*{1.7}}
    \put(29,46){\line(1,-1){2.5}}
    \put(27,14){\line(1,1){2.5}}
    \put(42,9){\line(0,1){3}}
    \put(20,50){$P_6$}
    \put(16,8){$P_1$}
    \put(38,1){$P_2$}
    \put(84,51){\line(0,-1){3}}
    \put(105,30){\line(-1,0){3}}
    \put(72,13){\line(1,1){2.5}}
    \put(80,53){$P_5$}
    \put(108,27){$P_4$}
    \put(65,6){$P_3$}
    \put(34,29){\circle*{1.7}}
    \put(49,32){\circle*{1.7}}
    \put(76,36){\circle*{1.7}}
    \put(94,29){\circle*{1.7}}
    \put(34,30){\line(0,1){8}}
    \put(49,33){\line(1,2){4}}
    \put(76,37){\line(1,2){3}}
    \put(94,30){\line(0,1){8}}
    \qbezier[7](34,28)(34,24)(34,21)
    \qbezier[7](49,31)(49,27)(49,24)
    \qbezier[7](76,35)(76,31)(76,28)
    \qbezier[7](94,28)(94,24)(94,21)
    \put(29,39){$Q_1$}
    \put(44,42){$Q_2$}
    \put(78,44){$Q_3$}
    \put(96,39){$Q_4$}
    \put(86,18){\circle*{1.7}}
    \put(88,17){\vector(1,-1){6}}
    \put(82,11){$ev_0$}
  \end{picture}
  \par\smallskip
  \renewcommand{\thefigure}{3.8.7}
  \refstepcounter{figure}
  \textbf{\figurename~\thefigure}\label{fig:3-8-7}
\end{figure}

式 \eqref{eq:3-8-79-4} 中出现
$GW_{0,\ell+1}(M)(y\otimes L)$ 的原因与 $\mathfrak p$ 的情形相同。
命题 3.8.78 的证明细节是本节已讨论论证的组合；第 7.4.11 小节将作
详细说明。

还要指出，第 7.4.10--7.4.11 小节将用 $L_\infty$ 结构重新表述
命题 3.8.78；见定理 7.4.192。

\begin{example}{3.8.80}\label{ex:3-8-80}
考虑
\[
 (M,\omega_M)
 = (\mathbb CP^1,\omega)\times(\mathbb CP^1,-\omega),
\]
并令 $L$ 为对角线，其定向选得使第一因子上的投影保持定向。取
$p=(z_0,w_0)\in M$，并考虑
$\operatorname{PD}[p]\in H^{2n}(M;\mathbb Q)$。在
\eqref{eq:3-8-79-4} 中置 $y=\operatorname{PD}[p]$，得到
\begin{equation}
\begin{aligned}
 &\delta_M\bigl(
 \mathfrak p_{1,0}(\operatorname{PD}[p]\otimes1)
 \bigr)
 +\mathfrak p_{1,1}\bigl(
 \operatorname{PD}[p]\otimes\mathfrak m_0(1)
 \bigr)\\
 &\quad+
 \mathfrak p_{0,1}\bigl(
 1\otimes\mathfrak q_{1,0}(\operatorname{PD}[p]\otimes1)
 \bigr)
 +GW_{0,2}(M)(\operatorname{PD}[p]\otimes[L])=0.
\end{aligned}
\tag{3.8.81}\label{eq:3-8-81}
\end{equation}
由维数计算，$\mathfrak m_0(1)=0$ 是零链。因此
\eqref{eq:3-8-81} 蕴含
\begin{equation}
\begin{aligned}
 &\left[
 \mathfrak p_{0,1}\bigl(
 1\otimes\mathfrak q_{1,0}(\operatorname{PD}[p]\otimes1)
 \bigr)
 \right]\\
 &\qquad=-
 \bigl[GW_{0,2}(M)(\operatorname{PD}[p]\otimes[L])\bigr].
\end{aligned}
\tag{3.8.82}\label{eq:3-8-82}
\end{equation}
下面直接检验式 \eqref{eq:3-8-82} 的首项。

写
$H^2(M;\mathbb Z)=\mathbb Z\oplus\mathbb Z$，并在 $M$ 上取复结构
\[
 J_M=J_{\mathbb CP^1}\oplus(-J_{\mathbb CP^1}).
\]
于是 $\widetilde\beta_1=(1,0)$ 与
$\widetilde\beta_2=(0,-1)$ 都由辛面积最小的伪全纯球面表示。我们有
\[
 \mathcal M_{0,3}(M;\widetilde\beta_1)
 \mathbin{\times_M}(\{p\}\times L)
 \cong S^2,
\]
其评价映射之像为
\[
 ev_0\left(
 \mathcal M_{0,3}(M;\widetilde\beta_1)
 \mathbin{\times_M}(\{p\}\times L)
 \right)
 =\{(z,w_0)\mid z\in\mathbb CP^1\}.
\]
因此
\begin{equation}
 GW_{0,2}(M;\widetilde\beta_1)
 (\operatorname{PD}[p]\otimes[L])
 =\mathbb CP^1\times\{w_0\}.
\tag{3.8.83.1}\label{eq:3-8-83-1}
\end{equation}
另一方面，计及
$\widetilde\beta_1\cap L=-\widetilde\beta_2\cap L$，得到
\begin{equation}
 GW_{0,2}(M;\widetilde\beta_2)
 (\operatorname{PD}[p]\otimes[L])
 =\{z_0\}\times\mathbb CP^1.
\tag{3.8.83.2}\label{eq:3-8-83-2}
\end{equation}
这里，\eqref{eq:3-8-83-2} 右端 $\mathbb CP^1$ 的定向与表示
$\widetilde\beta_2$ 的伪全纯球面的定向相反。因此
\begin{equation}
 \bigl[GW_{0,2}(M)(\operatorname{PD}[p]\otimes[L])\bigr]
 \equiv T^{\omega[\mathbb CP^1]}(1,1)
 \pmod{T^{\omega[\mathbb CP^1]}\Lambda^+_{0,\mathrm{nov}}}.
\tag{3.8.84}\label{eq:3-8-84}
\end{equation}

再来计算 \eqref{eq:3-8-82} 的左端。令
$\beta\in H_2(M,L;\mathbb Z)$ 为 $\widetilde\beta_i$ 的像；
$\widetilde\beta_1$ 与 $\widetilde\beta_2$ 在
$H_2(M,L;\mathbb Z)$ 中给出同一个元素。作认同
$\mathbb CP^1=\mathbb C\cup\{\infty\}$，并取
$z_0=0$、$w_0=\infty$。于是
\begin{equation}
 \mathcal M_{0,1}^{\mathrm{main,reg}}(L;\beta)_{ev}
 \mathbin{\times_M}\{p\}
 \cong(0,\infty),
\tag{3.8.85}\label{eq:3-8-85}
\end{equation}
且 $r\in(0,\infty)$ 对应于全纯映射
$u_r\colon D^2\to M$，其中
\[
 u_r(z)=\left(rz,\frac r{\overline z}\right).
\]
注意 $u_r(0)=p$ 且 $u_r(\partial D^2)\subseteq L$；见图 3.8.8。
由此可得
\[
 \mathcal M_{1,1}^{\mathrm{reg}}(L;\beta)_{ev}
 \mathbin{\times_M}\{p\}
 \cong\mathbb R\times S^1
\]
以及
\[
 \left[
 ev_0\left(
 \mathcal M_{1,1}^{\mathrm{main}}(L;\beta)_{ev}
 \mathbin{\times_M}\{p\}
 \right)
 \right]=\pm[L].
\]
也就是说，
\[
 \mathfrak q_{1,0}(\operatorname{PD}[p]\otimes1)
 \equiv\pm T^{\omega[\mathbb CP^1]}[L]
 \pmod{T^{\omega[\mathbb CP^1]}\Lambda^+_{0,\mathrm{nov}}}.
\]
所以
\begin{equation}
\begin{aligned}
 &\left[
 \mathfrak p_{0,1}\bigl(
 1\otimes\mathfrak q_{1,0}(\operatorname{PD}[p]\otimes1)
 \bigr)
 \right]\\
 &\qquad\equiv
 \pm T^{\omega[\mathbb CP^1]}[L]
 \pmod{T^{\omega[\mathbb CP^1]}\Lambda^+_{0,\mathrm{nov}}}.
\end{aligned}
\tag{3.8.86}\label{eq:3-8-86}
\end{equation}
这与 \eqref{eq:3-8-82}、\eqref{eq:3-8-84} 相容。（这里不试图核对
符号。）

还要指出，模空间 \eqref{eq:3-8-85} 的两个端部分别对应于常值圆盘
加上一个球面气泡；这两个气泡分别对应 \eqref{eq:3-8-83-1} 与
\eqref{eq:3-8-83-2}。

\begin{figure}[htbp]
  \centering
  \setlength{\unitlength}{0.65mm}
  \begin{picture}(108,78)(0,-1)
    \qbezier(24,65)(57,78)(91,62)
    \qbezier(91,62)(62,33)(30,9)
    \qbezier(30,9)(12,37)(24,65)
    \qbezier(24,65)(49,59)(91,62)
    \qbezier(24,65)(24,33)(30,9)
    \qbezier(24,65)(37,49)(50,33)
    \qbezier(24,65)(42,48)(60,31)
    \qbezier(24,65)(47,49)(70,33)
    \qbezier(24,65)(52,51)(80,37)
    \qbezier(24,65)(58,64)(91,62)
    \put(57,66){
      \begin{picture}(0,0)
        \qbezier(-4,0)(-4,8)(0,8)
        \qbezier(0,8)(4,8)(4,0)
        \qbezier(4,0)(4,-8)(0,-8)
        \qbezier(0,-8)(-4,-8)(-4,0)
      \end{picture}}
    \put(23,40){
      \begin{picture}(0,0)
        \qbezier(-8,0)(-8,4)(0,4)
        \qbezier(0,4)(8,4)(8,0)
        \qbezier(8,0)(8,-4)(0,-4)
        \qbezier(0,-4)(-8,-4)(-8,0)
      \end{picture}}
    \put(25,29){
      \begin{picture}(0,0)
        \qbezier(-8,0)(-8,4)(0,4)
        \qbezier(0,4)(8,4)(8,0)
        \qbezier(8,0)(8,-4)(0,-4)
        \qbezier(0,-4)(-8,-4)(-8,0)
      \end{picture}}
    \put(52,27){
      \begin{picture}(0,0)
        \qbezier(-4,0)(-4,8)(0,8)
        \qbezier(0,8)(4,8)(4,0)
        \qbezier(4,0)(4,-8)(0,-8)
        \qbezier(0,-8)(-4,-8)(-4,0)
      \end{picture}}
    \put(62,35){
      \begin{picture}(0,0)
        \qbezier(-4,0)(-4,8)(0,8)
        \qbezier(0,8)(4,8)(4,0)
        \qbezier(4,0)(4,-8)(0,-8)
        \qbezier(0,-8)(-4,-8)(-4,0)
      \end{picture}}
    \put(72,43){
      \begin{picture}(0,0)
        \qbezier(-4,0)(-4,8)(0,8)
        \qbezier(0,8)(4,8)(4,0)
        \qbezier(4,0)(4,-8)(0,-8)
        \qbezier(0,-8)(-4,-8)(-4,0)
      \end{picture}}
    \put(18,68){$p$}
    \put(86,29){$L$}
  \end{picture}
  \par\smallskip
  \renewcommand{\thefigure}{3.8.8}
  \refstepcounter{figure}
  \textbf{\figurename~\thefigure}\label{fig:3-8-8}
\end{figure}
\end{example}

现在令
\[
 \mathfrak b\in C^2(M;\Lambda^+_{0,\mathrm{nov}}),
 \qquad \delta_M\mathfrak b=0.
\]
定义
\[
 \mathfrak p_k^{\mathfrak b}\colon
 B_k^{\mathrm{cyc}}C_{1+}(L;\Lambda_{0,\mathrm{nov}})[1]
 \longrightarrow S^*(M;\Lambda_{0,\mathrm{nov}})
\]
如下：
\[
 \mathfrak p_k^{\mathfrak b}(x)
 =\sum_\ell
 \mathfrak p_{\ell,k}(\mathfrak b^\ell\otimes x)
 =\mathfrak p(e^{\mathfrak b}\otimes x).
\]
于是有
\begin{align}
 \mathfrak p_1^{\mathfrak b}
 &\equiv i_!\pmod{\Lambda^+_{0,\mathrm{nov}}},
 \tag{3.8.87.1}\label{eq:3-8-87-1}\\
 \mathfrak p^{\mathfrak b}\circ
 \widehat d^{\mathfrak b,\mathrm{cyc}}
 +\delta_M\circ\mathfrak p_k^{\mathfrak b}
 &=0
 \qquad(k>0),
 \tag{3.8.87.2}\label{eq:3-8-87-2}\\
 \mathfrak p_1^{\mathfrak b}\circ
 \mathfrak m_0^{\mathfrak b}(1)
 +\delta_M\circ\mathfrak p_0^{\mathfrak b}(1)
 +\sum_\ell\frac1{\ell!}
 GW_{0,\ell+1}(M)
 (\mathfrak b^{\otimes\ell}\otimes L)
 &=0.
 \tag{3.8.87.3}\label{eq:3-8-87-3}
\end{align}
它们与 (3.8.10.1)--(3.8.10.3) 是同样的公式。换言之，在研究
\[
 HF\bigl(
 (L;(\mathfrak b,b)),(L;(\mathfrak b,b'));
 \Lambda_{0,\mathrm{nov}}
 \bigr)
\]
时，可以用 $\mathfrak p_k^{\mathfrak b}$ 代替 $\mathfrak p_k$。

下面推广定理 3.8.11。回顾定义 3.8.40 中
\[
 \widehat{\mathcal M}^{\mathfrak b}(L)
 =\left\{
 b\in C(L;\Lambda^+_{0,\mathrm{nov}})
 \ \middle|\
 (\mathfrak b,b)\in\widehat{\mathcal M}_{\mathrm{def}}(L)
 \right\},
\]
而 $\widehat{\mathcal M}^{\mathfrak b}_{\mathrm{weak}}(L)$ 以类似方式定义。
以 $[o_i(L;\mathfrak b)]$ 表示
$\widehat{\mathcal M}^{\mathfrak b}(L)$ 非空的障碍类（见定理 3.6.18）；
$[o_i(L;\mathfrak b;\mathrm{weak})]$ 以类似方式定义（见定理 3.6.43）。

\begin{theorem}{3.8.88}\label{thm:3-8-88}
设 $L\subset M$ 是相对 spin 拉格朗日子流形，并设
$\mathfrak b\in H^2(M;\Lambda^+_{0,\mathrm{nov}})$。

\smallskip
\noindent\textup{(3.8.89.1)}
若
\begin{equation}
 \sum_{\ell=0}^{\infty}\frac1{\ell!}
 \left[
 GW_{0,\ell+1}(M)
 (\mathfrak b^{\otimes\ell}\otimes L)
 \right]\ne0,
\tag{3.8.90}\label{eq:3-8-90}
\end{equation}
则 $\widehat{\mathcal M}^{\mathfrak b}(L)=\varnothing$。
\liftIndex{I138}

\noindent\textup{(3.8.89.2)}
若 \eqref{eq:3-8-90} 左端为零，则障碍类
$[o_i(L;\mathfrak b)]$ 属于 Gysin 同态
\[
 i_!\colon H^{2-\mu_i}(L;\mathbb Q)
 \longrightarrow H^{n+2-\mu_i}(M;\mathbb Q)
\]
的核。

\noindent\textup{(3.8.89.3)}
若 \eqref{eq:3-8-90} 左端不具有
$c\operatorname{PD}_M[L]$ 的形式，则
$\widehat{\mathcal M}^{\mathfrak b}_{\mathrm{weak}}(L)=\varnothing$。

\noindent\textup{(3.8.89.4)}
若 \eqref{eq:3-8-90} 左端具有
$c\operatorname{PD}_M[L]$ 的形式，则\translatornote{原书此项把
$o_i(L;\mathfrak b;\mathrm{weak})$ 直接写作 Gysin 同态之核中的元素，
推论 3.8.96 (2) 又把其代表元直接送入上同调上的 $i_!$；两处都漏去了
上同调类方括号。同一定理的 (3.8.89.2)、障碍类定义以及推论的假设均
明确使用 $[o_i]$，译文据此统一补回方括号。}
\[
 [o_i(L;\mathfrak b;\mathrm{weak})]
 \in\operatorname{Ker}\left(
 i_!\colon H^{2-\mu_i}(L;\mathbb Q)
 \longrightarrow H^{n+2-\mu_i}(M;\mathbb Q)
 \right).
\]

\noindent\textup{(3.8.89.5)}
若 \eqref{eq:3-8-90} 左端等于
$c\operatorname{PD}_M[L]$，且
$b\in\widehat{\mathcal M}^{\mathfrak b}_{\mathrm{weak}}(L)$、
$[L]\ne0$ 于 $H(M;\mathbb Q)$ 中，则
\[
 c=\mathfrak{PO}(\mathfrak b,b).
\]
\end{theorem}

定理 3.8.88 的证明与使用 (3.8.87) 的定理 3.8.11 之证明相同，故略去。

\begin{remark}{3.8.91}\label{rem:3-8-91}
这里不能把引理 3.8.14 直接推广到当前情形。事实上，
\eqref{eq:3-8-90} 的左端在 $*=2n$ 或 $2n-1$ 时可能具有非零的
$H^*(M;\mathbb Q)\otimes\Lambda^+_{0,\mathrm{nov}}$ 分量。

若
$\mathfrak b\in H^n(M;\mathbb Q)\otimes\Lambda^+_{0,\mathrm{nov}}$ 且
$\mathfrak b\cap L\ne0$，则
\[
 GW_{0,2}(M;\widetilde\beta_0)(\mathfrak b,L)
 =\operatorname{PD}_M(\mathfrak b\cap L)\ne0.
\]
这里 $\widetilde\beta_0=0$，而
\[
 \operatorname{PD}_M\colon
 H_d(M;\mathbb Q)\otimes\Lambda_{0,\mathrm{nov}}
 \longrightarrow
 H^{2n-d}(M;\mathbb Q)\otimes\Lambda_{0,\mathrm{nov}}
\]
是 $M$ 上的庞加莱对偶。
\end{remark}

引理 3.8.14 可以作如下部分推广。

\begin{lemma}{3.8.92}\label{lem:3-8-92}
设 $B_j\in H(M;\mathbb Q)$（$j=1,\ldots,\ell$），并假设
\[
 GW_{0,\ell+1}(M;\widetilde\beta)
 (B_1,\ldots,B_\ell,L)
 \in H^*(M;\mathbb Q)\otimes\Lambda^+_{0,\mathrm{nov}},
\]
其中 $*=2n$ 或 $2n-1$。再假设下列条件之一成立：
\begin{align}
 \widetilde\beta&\ne0
 \quad\text{于 }H_2(M;\mathbb Z),
 \tag{3.8.93.1}\label{eq:3-8-93-1}\\
 \ell&\ne1.
 \tag{3.8.93.2}\label{eq:3-8-93-2}
\end{align}
则
\[
 \left[
 GW_{0,\ell+1}(M;\widetilde\beta)
 (B_1,\ldots,B_\ell,L)
 \right]=0.
\]
\end{lemma}

\begin{proof}
置 $d_j=\deg B_j=2n-\dim B_j$。有
\[
 *=n-2c_1(M)[\widetilde\beta]+2
 +\sum_{j=1}^{\ell}(d_j-2),
\]
这是因为
\[
 \dim\mathcal M_{0,\ell+2}(M;\widetilde\beta)
 =2n+2c_1(M)[\widetilde\beta]-2+2\ell.
\]
由此，模空间
\begin{equation}
 \mathcal M_{0,\ell+1}(M;\widetilde\beta)
 \mathbin{\times_{M^\ell}}
 (B_1\times\cdots\times B_\ell)
\tag{3.8.94}\label{eq:3-8-94}
\end{equation}
的虚维数为
\[
 2n-4+2c_1(M)[\widetilde\beta]
 +\sum_{j=1}^{\ell}(2-d_j)
 =3n-*-2.
\]
所以由假设，
\[
 \dim\left(
 \mathcal M_{0,\ell+1}(M;\widetilde\beta)
 \mathbin{\times_{M^\ell}}
 (B_1\times\cdots\times B_\ell)
 \right)
 =n-2\quad\text{或}\quad n-1.
\]
因此纤维积
\[
 \mathcal M_{0,\ell+1}(M;\widetilde\beta)
 \mathbin{\times_{M^{\ell+1}}}
 (B_1\times\cdots\times B_\ell\times L)
\]
为空。可以在
\begin{equation}
 \mathcal M_{0,\ell+2}(M;\widetilde\beta)
 \mathbin{\times_{M^{\ell+1}}}
 (B_1\times\cdots\times B_\ell\times L)
\tag{3.8.95}\label{eq:3-8-95}
\end{equation}
上选取与遗忘映射
\[
\begin{aligned}
 &\mathcal M_{0,\ell+2}(M;\widetilde\beta)
 \mathbin{\times_{M^{\ell+1}}}
 (B_1\times\cdots\times B_\ell\times L)\\
 &\quad\longrightarrow
 \mathcal M_{0,\ell+1}(M;\widetilde\beta)
 \mathbin{\times_{M^{\ell+1}}}
 (B_1\times\cdots\times B_\ell\times L)
\end{aligned}
\]
相容的扰动，并使 $ev_0$ 与该遗忘映射相容。
\translatornote{原书与引理 3.8.14 的证明一样，先称该模空间间的映射为
forgetful map，下一分句却说 $ev_0$ 与 evaluation map 相容。此映射的
源、靶均为模空间，作用是忘掉一个标记点；译文沿用引理 3.8.14 已作的
裁决，统一称为“遗忘映射”。}
见 \textbf{[FuOn99II]} 第 23 节。这里用到了假设 (3.8.93)：事实上，
当 $\widetilde\beta=\widetilde\beta_0=0$ 且 $\ell=1$ 时，带两个标记点的
常值球面映射不稳定，故
$\mathcal M_{0,2}(M;\widetilde\beta_0)$ 为空；另一方面，
$\mathcal M_{0,3}(M;\widetilde\beta_0)=M$，因而非空。

在给定假设下，\eqref{eq:3-8-95} 的相容扰动给出空的虚基本链，所以
\[
 \left[
 GW_{0,\ell+1}(M;\widetilde\beta)
 (B_1,\ldots,B_\ell,L)
 \right]=0.
\]
\end{proof}

令
\[
 \Pi^d\colon H(M;\Lambda_{0,\mathrm{nov}})
 \longrightarrow H^d(M;\mathbb Q)\otimes\Lambda_{0,\mathrm{nov}}
\]
为投影。由引理 3.8.92 与 \eqref{eq:3-8-87-3}，当
$*=2n-1$ 或 $2n$ 时有
\[
 \Pi^*\left(
 \mathfrak p_1^{\mathfrak b}\circ\mathfrak m_0^{\mathfrak b}(1)
 +\delta_M\circ\mathfrak p_0^{\mathfrak b}(1)
 \right)
 =\Pi^*\bigl(
 \operatorname{PD}_M([L]\cap\mathfrak b)
 \bigr).
\]
利用此式可以证明下面推广推论 3.8.16 的结果。

\begin{corollary}{3.8.96}\label{cor:3-8-96}
假设下列两个条件之一成立：
\begin{align*}
 \textup{(a)}\quad
 \Pi^n\bigl(\operatorname{PD}_L(\mathfrak b\cap L)\bigr)&=0,
 &\text{或}\qquad
 \textup{(b)}\quad
 \Pi^{n-1}\bigl(\operatorname{PD}_L(\mathfrak b\cap L)\bigr)&=0.
\end{align*}
这里
\[
 \operatorname{PD}_L\colon
 H_d(L;\mathbb Q)\otimes\Lambda_{0,\mathrm{nov}}
 \longrightarrow
 H^{n-d}(L;\mathbb Q)\otimes\Lambda_{0,\mathrm{nov}}
\]
是 $L$ 上的庞加莱对偶。

\begin{enumerate}[label=\textup{(\arabic*)},leftmargin=2.8em]
\item 假设 (a) 成立，且障碍类 $[o_i(L;\mathfrak b)]$ 已定义并属于
$H^n(L;\mathbb Q)$，则 $[o_i(L;\mathfrak b)]=0$。

\item 假设 (b) 成立。若
$[o_i(L;\mathfrak b)]\in H^{n-1}(L;\mathbb Q)$，则
\[
 i_!\bigl([o_i(L;\mathfrak b)]\bigr)=0.
\]
\end{enumerate}
同样的结论也适用于
$[o_i(L;\mathfrak b;\mathrm{weak})]$。
\end{corollary}

利用推论 3.8.96，可以证明下面两个部分推广推论 3.8.19 的结果。

\begin{corollary}{3.8.97}\label{cor:3-8-97}
设 $L$ 是作为有理同调球面的相对 spin 拉格朗日子流形，并设
\[
 \mathfrak b\in
 \bigoplus_{*\geq2}H^*(M;\mathbb Q)
 \otimes\Lambda^+_{0,\mathrm{nov}}.
\]
假设
$\Pi^n(\operatorname{PD}_L(\mathfrak b\cap L))=0$，并且不存在边界
落在 $L$ 上、马斯洛夫指标非负的伪全纯圆盘。则
\begin{equation}
 \sum_{\ell=0}^{\infty}\frac1{\ell!}
 \left[
 GW_{0,\ell+1}(M)
 (\mathfrak b^{\otimes\ell}\otimes L)
 \right]=0.
\tag{3.8.98}\label{eq:3-8-98}
\end{equation}
\end{corollary}

\begin{proof}
由于 $L$ 是有理同调球面，$[o_i(L;\mathfrak b)]$ 或属于
$H^n(L;\mathbb Q)$，或属于 $H^0(L;\mathbb Q)$。上述伪全纯圆盘不存在，
且 $\mathfrak b$ 的次数不小于 $2$，故\translatornote{原书下式把
$\mathfrak m_0^{\mathfrak b}(1)$ 写成 $M$ 的上同调中的元素；然而
$\mathfrak m^{\mathfrak b}$ 是 $L$ 上的 bulk 形变 $A_\infty$ 结构，
其零元运算的输出属于 $L$ 的复形。译文依算子型别把 $H^*(M)$ 修正为
$H^*(L)$。}
\[
 \mathfrak m_0^{\mathfrak b}(1)
 \in\bigoplus_{*\geq2}H^*(L;\mathbb Q)
 \otimes\Lambda^+_{0,\mathrm{nov}}.
\]
所以 $[o_i(L;\mathfrak b)]$ 不属于 $H^0(L;\mathbb Q)$。推论 3.8.96
因而给出 $[o_i(L;\mathfrak b)]=0$，再由定理 3.8.88 (1) 即得本推论。
\end{proof}

\begin{corollary}{3.8.99}\label{cor:3-8-99}
设 $L$ 是马斯洛夫指标恒为零的相对 spin 拉格朗日子流形，并设
$\mathfrak b\in H^2(M;\mathbb Q)\otimes\Lambda^+_{0,\mathrm{nov}}$。
若 $H^2(L;\mathbb Q)=0$，则 \eqref{eq:3-8-98} 成立。
\end{corollary}

\begin{proof}
由次数计算可知，$[o_i(L;\mathfrak b)]$ 属于
$H^2(L;\mathbb Q)$。
\translatornote{原书这里写成 $H^2(M;\mathbb Q)$；但
$o_i(L;\mathfrak b)$ 是 $L$ 上的障碍上链，且本推论正以
$H^2(L;\mathbb Q)=0$ 推出该类消失。译文依对象型别及紧接着的推论逻辑
恢复为 $H^2(L;\mathbb Q)$。}
故由假设，$\widehat{\mathcal M}^{\mathfrak b}(L)\ne\varnothing$。
定理 3.8.88 (1) 于是给出本推论。
\end{proof}
